\chapter{Introducción}

\section{Requisitos Previos}
\begin{itemize}
    \item Servidor web Apache instalado y configurado.
    \item Certificados de cliente necesarios: certificado de la autoridad de la asignatura y un certificado secundario (ej. FNMT).
    \item Herramientas OpenSSL disponibles.
\end{itemize}

\section{Objetivos de la práctica}
\begin{itemize}
    \item Modificar la configuración del servidor web Apache para que acepte certificados de cliente firmados por la misma autoridad de certificación y por una autoridad adicional (preferentemente FNMT).
    \item Habilitar la ejecución de scripts CGI/binarios (bash) o Python para interactuar con los datos del certificado.
    \item Instalar un script que reemplace el fichero \texttt{index.html} para mostrar los datos más importantes del certificado presentado por el cliente (Common Name, Issuer, Validity, etc.) a través de variables de entorno.
    \item Demostrar la capacidad del servidor web para aceptar dos conexiones distintas utilizando certificados firmados por autoridades diferentes.
    \item Verificar la correcta comunicación y aceptación de certificados mediante comandos \texttt{openssl s\_client} y navegadores (Firefox), incluyendo la visualización de los datos del certificado en el script.
\end{itemize}

\section{Fundamentos teóricos}
Para la realización de esta práctica, se aplican diversos conceptos teóricos clave de la ciberseguridad y la gestión de redes:
\begin{itemize}
    \item Certificados Digitales y Autoridades de Certificación (CA): Entender cómo funcionan los certificados X.509, la estructura de una cadena de confianza y el rol de las Autoridades de Certificación (como FNMT) en la emisión y validación de identidades digitales.
    \item Configuración del Servidor Web SSL/TLS: Comprender cómo configurar un servidor web (Apache) para que gestione la autenticación de clientes mediante certificados (usando módulos como \texttt{mod\_ssl}) y cómo gestionar múltiples CAs confiables.
    \item Validación de Certificados en el Handshake TLS: Analizar el proceso de negociación TLS, donde el cliente presenta su certificado al servidor. Esto implica verificar la validez del certificado del cliente y asegurar que la cadena de confianza (la ruta de la CA) es válida para establecer una conexión segura.
    \item Interacción con el Sistema Operativo (Scripts CGI/Python): Entender cómo los scripts pueden interactuar con el sistema operativo y extraer metadatos específicos (como Common Name, Issuer, Validity) contenidos dentro del certificado digital para presentarlos al usuario o al servidor web.
    \item Mecanismos de Autenticación Múltiple: Demostrar la capacidad del servidor para validar certificados emitidos por diferentes CAs simultáneamente, lo que requiere una configuración adecuada en el servidor para manejar múltiples \textit{trust stores}.
    \item Herramientas de Diagnóstico (OpenSSL y Navegadores): Utilizar herramientas como \texttt{openssl s\_client} para verificar directamente la comunicación SSL/TLS, la presencia de certificados y la correcta validación de las firmas entre cliente y servidor.
\end{itemize}

\section{Enunciado de la práctica}
El desarrollo del ejercicio y de la presente memoria seguirán el guion de actividades propuesto en el enunciado:

\begin{enumerate}
    \item \textbf{Modificación de la configuración del Apache}
    \begin{itemize}
        \item Configurar el servidor para que solicite certificados de cliente y acepte los firmados por la misma autoridad de la asignatura y por otra autoridad (preferentemente FNMT).
        \item Habilitar la ejecución de scripts CGI-bin (\texttt{bash}), código Python, etc.
    \end{itemize}
    \item \textbf{Instalación del script de sustitución de index.html}
    \begin{itemize}
        \item Breve explicación del script utilizado y visualización de sus permisos.
    \end{itemize}
    \item \textbf{Carga del segundo certificado de cliente en el Firefox}
    \begin{itemize}
        \item Preferentemente, utilizar un certificado FNMT real. Alternativamente, crear una autoridad de certificación nueva, generar un segundo certificado de cliente para el usuario con esta nueva autoridad e instalarlo en el Firefox, junto con el certificado de la autoridad.
        \item Si se utiliza el certificado FNMT, los certificados de autoridad ya se encuentran en el directorio del sistema Linux (o no) y deben incorporarse a la configuración del Apache.
        \item Si se utiliza otra autoridad, repetir los pasos para su certificado raíz y asegurarse de que se utiliza la extensión \texttt{-addext "basicConstraints=CA:TRUE"} en el comando de creación del certificado autofirmado de autoridad.
    \end{itemize}
    \item \textbf{Rearranque del Apache y verificación del fichero error\_log}
    \begin{itemize}
        \item Realizar el reinicio del servicio Apache.
        \item Verificar el contenido del fichero \texttt{error\_log}.
    \end{itemize}
    \item \textbf{Conexión de prueba con el comando “openssl s\_client” libre de errores}
    \begin{itemize}
        \item Ejecutar el siguiente comando para verificar la conexión y los certificados:
        \[
        \texttt{openssl s\_client www.ejemplo.com:443 -cert certcliente.pem -key certcliente.key -CAfile cert\_raiz.pem}
        \]
        \item Verificar que no se producen errores en la conexión ni en los certificados.
        \item Verificar que NO hay errores en \texttt{error\_log}.
    \end{itemize}
    \item \textbf{Conexión sin errores desde Firefox con el servidor, con el certificado de cliente de la autoridad de la asignatura}
    \begin{itemize}
        \item Asegurar que el script muestre los datos más importantes del certificado de cliente.
        \item Mostrar captura de pantalla de la conexión (ejecución del script).
    \end{itemize}
    \item \textbf{Conexión sin errores desde Firefox con el servidor, con el certificado de cliente FNMT o uno secundario creado por el estudiante}
    \begin{itemize}
        \item Asegurar que el script muestre los datos más importantes del certificado de cliente.
        \item Mostrar captura de pantalla de la conexión (ejecución del script).
    \end{itemize}
\end{enumerate}